\clearpage\phantomsection\label{couverture-programme}\pdfbookmark[0]{Repérage des attendus}{bm-couverture-programme}
\gdef\otChapNum{0}\pagestyle{otFront}\otPageTitre{Repérage des attendus}{Troisième 2026-2027 · Lecture par notions}
Cette matrice rassemble les attendus annuels, avec les acquis du cycle 4
réactivés pour les utiliser. Les quatre repères de chaque chapitre renvoient
à des exercices précis ; les séries citées ne sont pas exhaustives.
Les dix-huit défis prolongent ces notions sans ajouter d’exigence à l’examen.\par\bigskip
Points à retrouver : puissances négatives (2), différence de carrés et
produits nuls (4-5), histogrammes de même amplitude (10), arbres à deux
étapes (11), Thalès en papillon (13), homothéties négatives (15), sections,
assemblages et repérage sur la sphère (16), scripts parallèles (17).\par\bigskip
\begin{attention}La précision perdue par un regroupement statistique en classes est signalée :
les valeurs extrêmes exactes et la médiane exacte ne se déduisent pas toujours
d’un histogramme. La programmation est abordée par blocs, comme au collège.\end{attention}
\clearpage
\otPageTitre{Repères et exercices}{Chapitres 1 à 3}\fontsize{10}{13}\selectfont
\otIntertitre{1. Nombres rationnels et fractions}\begin{tabularx}{\linewidth}{@{}Yl@{}}\toprule
Priorités et signes & \hyperref[exo-1-1]{1.1}, \hyperref[exo-1-2]{1.2}\\[2mm]
Opérations sur les fractions & \hyperref[exo-1-3]{1.3}, \hyperref[exo-1-4]{1.4}, \hyperref[exo-1-5]{1.5}, \hyperref[exo-1-7]{1.7}, \hyperref[exo-1-8]{1.8}\\[2mm]
Comparaison et arrondis & \hyperref[exo-1-6]{1.6}, \hyperref[exo-1-10]{1.10}, \hyperref[exo-1-15]{1.15}\\[2mm]
Fractions dans les problèmes & \hyperref[exo-1-9]{1.9}, \hyperref[exo-1-11]{1.11}, \hyperref[exo-1-12]{1.12}, \hyperref[exo-1-13]{1.13}, \hyperref[exo-1-14]{1.14}\\[2mm]
\bottomrule\end{tabularx}\par\vspace{4mm}
\otIntertitre{2. Puissances, racines et écriture scientifique}\begin{tabularx}{\linewidth}{@{}Yl@{}}\toprule
Puissances entières et signes & \hyperref[exo-2-1]{2.1}, \hyperref[exo-2-2]{2.2}, \hyperref[exo-2-3]{2.3}, \hyperref[exo-2-11]{2.11}\\[2mm]
Notation scientifique & \hyperref[exo-2-4]{2.4}, \hyperref[exo-2-6]{2.6}, \hyperref[exo-2-7]{2.7}, \hyperref[exo-2-12]{2.12}, \hyperref[exo-2-13]{2.13}\\[2mm]
Racines et encadrements & \hyperref[exo-2-5]{2.5}, \hyperref[exo-2-8]{2.8}, \hyperref[exo-2-9]{2.9}, \hyperref[exo-2-10]{2.10}, \hyperref[exo-2-14]{2.14}\\[2mm]
Grandeurs et ordres de grandeur & \hyperref[exo-2-11]{2.11}, \hyperref[exo-2-12]{2.12}, \hyperref[exo-2-13]{2.13}, \hyperref[exo-2-14]{2.14}, \hyperref[exo-2-15]{2.15}\\[2mm]
\bottomrule\end{tabularx}\par\vspace{4mm}
\otIntertitre{3. Arithmétique et nombres premiers}\begin{tabularx}{\linewidth}{@{}Yl@{}}\toprule
Divisibilité et nombres premiers & \hyperref[exo-3-1]{3.1}, \hyperref[exo-3-2]{3.2}, \hyperref[exo-3-5]{3.5}, \hyperref[exo-3-6]{3.6}, \hyperref[exo-3-7]{3.7}\\[2mm]
Décomposition et fractions irréductibles & \hyperref[exo-3-3]{3.3}, \hyperref[exo-3-4]{3.4}, \hyperref[exo-3-8]{3.8}, \hyperref[exo-3-9]{3.9}\\[2mm]
Diviseurs communs et partages & \hyperref[exo-3-10]{3.10}, \hyperref[exo-3-11]{3.11}, \hyperref[exo-3-13]{3.13}\\[2mm]
Multiples communs et périodicité & \hyperref[exo-3-12]{3.12}, \hyperref[exo-3-14]{3.14}, \hyperref[exo-3-15]{3.15}\\[2mm]
\bottomrule\end{tabularx}\par\vspace{4mm}
\clearpage
\otPageTitre{Repères et exercices}{Chapitres 4 à 6}\fontsize{10}{13}\selectfont
\otIntertitre{4. Développer, factoriser et prouver}\begin{tabularx}{\linewidth}{@{}Yl@{}}\toprule
Substitution, opposé et réduction & \hyperref[exo-4-1]{4.1}, \hyperref[exo-4-2]{4.2}, \hyperref[exo-4-3]{4.3}\\[2mm]
Simple et double distributivité & \hyperref[exo-4-4]{4.4}, \hyperref[exo-4-7]{4.7}\\[2mm]
Factorisation et différences de carrés & \hyperref[exo-4-5]{4.5}, \hyperref[exo-4-6]{4.6}, \hyperref[exo-4-8]{4.8}, \hyperref[exo-4-9]{4.9}\\[2mm]
Identités et preuves par le calcul & \hyperref[exo-4-10]{4.10}, \hyperref[exo-4-11]{4.11}, \hyperref[exo-4-12]{4.12}, \hyperref[exo-4-13]{4.13}, \hyperref[exo-4-14]{4.14}, \hyperref[exo-4-15]{4.15}\\[2mm]
\bottomrule\end{tabularx}\par\vspace{4mm}
\otIntertitre{5. Équations et mise en équation}\begin{tabularx}{\linewidth}{@{}Yl@{}}\toprule
Équations du premier degré & \hyperref[exo-5-1]{5.1}, \hyperref[exo-5-2]{5.2}, \hyperref[exo-5-3]{5.3}, \hyperref[exo-5-4]{5.4}, \hyperref[exo-5-5]{5.5}, \hyperref[exo-5-6]{5.6}\\[2mm]
Équations produits et carrés & \hyperref[exo-5-8]{5.8}, \hyperref[exo-5-9]{5.9}, \hyperref[exo-5-10]{5.10}, \hyperref[exo-5-15]{5.15}\\[2mm]
Cas particuliers et validation & \hyperref[exo-5-7]{5.7}, \hyperref[exo-5-10]{5.10}\\[2mm]
Modélisation de situations & \hyperref[exo-5-11]{5.11}, \hyperref[exo-5-12]{5.12}, \hyperref[exo-5-13]{5.13}, \hyperref[exo-5-14]{5.14}, \hyperref[exo-5-15]{5.15}\\[2mm]
\bottomrule\end{tabularx}\par\vspace{4mm}
\otIntertitre{6. Proportionnalité, ratios et pourcentages}\begin{tabularx}{\linewidth}{@{}Yl@{}}\toprule
Proportionnalité et coefficients & \hyperref[exo-6-1]{6.1}, \hyperref[exo-6-2]{6.2}, \hyperref[exo-6-10]{6.10}\\[2mm]
Ratios et échelles & \hyperref[exo-6-5]{6.5}, \hyperref[exo-6-8]{6.8}, \hyperref[exo-6-9]{6.9}\\[2mm]
Évolutions et prix initiaux & \hyperref[exo-6-3]{6.3}, \hyperref[exo-6-4]{6.4}, \hyperref[exo-6-6]{6.6}, \hyperref[exo-6-7]{6.7}, \hyperref[exo-6-14]{6.14}, \hyperref[exo-6-15]{6.15}\\[2mm]
Comparaisons et pourcentages emboîtés & \hyperref[exo-6-11]{6.11}, \hyperref[exo-6-12]{6.12}, \hyperref[exo-6-13]{6.13}\\[2mm]
\bottomrule\end{tabularx}\par\vspace{4mm}
\clearpage
\otPageTitre{Repères et exercices}{Chapitres 7 à 9}\fontsize{10}{13}\selectfont
\otIntertitre{7. Grandeurs composées et conversions}\begin{tabularx}{\linewidth}{@{}Yl@{}}\toprule
Conversions de longueurs, aires et volumes & \hyperref[exo-7-1]{7.1}, \hyperref[exo-7-2]{7.2}, \hyperref[exo-7-3]{7.3}\\[2mm]
Durées et vitesses moyennes & \hyperref[exo-7-4]{7.4}, \hyperref[exo-7-6]{7.6}, \hyperref[exo-7-7]{7.7}, \hyperref[exo-7-11]{7.11}\\[2mm]
Débits et masses volumiques & \hyperref[exo-7-5]{7.5}, \hyperref[exo-7-8]{7.8}, \hyperref[exo-7-10]{7.10}, \hyperref[exo-7-12]{7.12}\\[2mm]
Grandeurs composées et contrôle des unités & \hyperref[exo-7-9]{7.9}, \hyperref[exo-7-13]{7.13}, \hyperref[exo-7-14]{7.14}, \hyperref[exo-7-15]{7.15}\\[2mm]
\bottomrule\end{tabularx}\par\vspace{4mm}
\otIntertitre{8. Fonctions : images et antécédents}\begin{tabularx}{\linewidth}{@{}Yl@{}}\toprule
Vocabulaire et calcul d’images & \hyperref[exo-8-1]{8.1}, \hyperref[exo-8-3]{8.3}, \hyperref[exo-8-4]{8.4}, \hyperref[exo-8-6]{8.6}\\[2mm]
Tableaux, courbes et antécédents & \hyperref[exo-8-2]{8.2}, \hyperref[exo-8-5]{8.5}, \hyperref[exo-8-7]{8.7}, \hyperref[exo-8-8]{8.8}, \hyperref[exo-8-10]{8.10}, \hyperref[exo-8-15]{8.15}\\[2mm]
Comparaison de fonctions & \hyperref[exo-8-9]{8.9}\\[2mm]
Modèles et domaines de validité & \hyperref[exo-8-11]{8.11}, \hyperref[exo-8-12]{8.12}, \hyperref[exo-8-13]{8.13}, \hyperref[exo-8-14]{8.14}\\[2mm]
\bottomrule\end{tabularx}\par\vspace{4mm}
\otIntertitre{9. Fonctions linéaires et affines}\begin{tabularx}{\linewidth}{@{}Yl@{}}\toprule
Fonctions linéaires et proportionnalité & \hyperref[exo-9-1]{9.1}, \hyperref[exo-9-2]{9.2}, \hyperref[exo-9-4]{9.4}, \hyperref[exo-9-6]{9.6}\\[2mm]
Fonctions affines et tracés & \hyperref[exo-9-3]{9.3}, \hyperref[exo-9-5]{9.5}, \hyperref[exo-9-7]{9.7}, \hyperref[exo-9-8]{9.8}\\[2mm]
Pente et ordonnée à l’origine & \hyperref[exo-9-3]{9.3}, \hyperref[exo-9-7]{9.7}, \hyperref[exo-9-8]{9.8}, \hyperref[exo-9-9]{9.9}\\[2mm]
Comparaison de modèles et seuils & \hyperref[exo-9-10]{9.10}, \hyperref[exo-9-11]{9.11}, \hyperref[exo-9-12]{9.12}, \hyperref[exo-9-13]{9.13}, \hyperref[exo-9-14]{9.14}, \hyperref[exo-9-15]{9.15}\\[2mm]
\bottomrule\end{tabularx}\par\vspace{4mm}
\clearpage
\otPageTitre{Repères et exercices}{Chapitres 10 à 12}\fontsize{10}{13}\selectfont
\otIntertitre{10. Statistiques et représentations}\begin{tabularx}{\linewidth}{@{}Yl@{}}\toprule
Effectifs et fréquences & \hyperref[exo-10-1]{10.1}, \hyperref[exo-10-2]{10.2}\\[2mm]
Moyennes simples et pondérées & \hyperref[exo-10-3]{10.3}, \hyperref[exo-10-6]{10.6}, \hyperref[exo-10-10]{10.10}, \hyperref[exo-10-11]{10.11}, \hyperref[exo-10-12]{10.12}\\[2mm]
Médiane et étendue & \hyperref[exo-10-4]{10.4}, \hyperref[exo-10-5]{10.5}, \hyperref[exo-10-7]{10.7}, \hyperref[exo-10-13]{10.13}, \hyperref[exo-10-14]{10.14}\\[2mm]
Diagrammes et histogrammes de même amplitude & \hyperref[exo-10-8]{10.8}, \hyperref[exo-10-9]{10.9}, \hyperref[exo-10-15]{10.15}\\[2mm]
\bottomrule\end{tabularx}\par\vspace{4mm}
\otIntertitre{11. Probabilités et expériences à deux étapes}\begin{tabularx}{\linewidth}{@{}Yl@{}}\toprule
Équiprobabilité et événements contraires & \hyperref[exo-11-1]{11.1}, \hyperref[exo-11-2]{11.2}, \hyperref[exo-11-3]{11.3}, \hyperref[exo-11-5]{11.5}, \hyperref[exo-11-10]{11.10}\\[2mm]
Tableaux et arbres à deux étapes & \hyperref[exo-11-6]{11.6}, \hyperref[exo-11-7]{11.7}, \hyperref[exo-11-14]{11.14}\\[2mm]
Tirages avec et sans remise & \hyperref[exo-11-8]{11.8}, \hyperref[exo-11-9]{11.9}, \hyperref[exo-11-11]{11.11}, \hyperref[exo-11-13]{11.13}\\[2mm]
Fréquences et simulations & \hyperref[exo-11-4]{11.4}, \hyperref[exo-11-15]{11.15}\\[2mm]
\bottomrule\end{tabularx}\par\vspace{4mm}
\otIntertitre{12. Pythagore et distances}\begin{tabularx}{\linewidth}{@{}Yl@{}}\toprule
Calculer une longueur avec Pythagore & \hyperref[exo-12-1]{12.1}, \hyperref[exo-12-2]{12.2}, \hyperref[exo-12-3]{12.3}, \hyperref[exo-12-4]{12.4}\\[2mm]
Réciproque et contraposée & \hyperref[exo-12-5]{12.5}, \hyperref[exo-12-6]{12.6}, \hyperref[exo-12-10]{12.10}\\[2mm]
Distances dans un repère orthonormé & \hyperref[exo-12-9]{12.9}\\[2mm]
Problèmes de longueurs et d’aires & \hyperref[exo-12-7]{12.7}, \hyperref[exo-12-8]{12.8}, \hyperref[exo-12-11]{12.11}, \hyperref[exo-12-12]{12.12}, \hyperref[exo-12-13]{12.13}, \hyperref[exo-12-14]{12.14}, \hyperref[exo-12-15]{12.15}\\[2mm]
\bottomrule\end{tabularx}\par\vspace{4mm}
\clearpage
\otPageTitre{Repères et exercices}{Chapitres 13 à 15}\fontsize{10}{13}\selectfont
\otIntertitre{13. Thalès et triangles semblables}\begin{tabularx}{\linewidth}{@{}Yl@{}}\toprule
Thalès emboîté et en papillon & \hyperref[exo-13-1]{13.1}, \hyperref[exo-13-2]{13.2}, \hyperref[exo-13-3]{13.3}, \hyperref[exo-13-4]{13.4}, \hyperref[exo-13-5]{13.5}, \hyperref[exo-13-8]{13.8}, \hyperref[exo-13-9]{13.9}\\[2mm]
Réciproque et non-parallélisme & \hyperref[exo-13-6]{13.6}, \hyperref[exo-13-7]{13.7}, \hyperref[exo-13-14]{13.14}\\[2mm]
Triangles semblables et correspondances & \hyperref[exo-13-10]{13.10}, \hyperref[exo-13-12]{13.12}\\[2mm]
Rapports des longueurs, périmètres et aires & \hyperref[exo-13-11]{13.11}, \hyperref[exo-13-13]{13.13}, \hyperref[exo-13-15]{13.15}\\[2mm]
\bottomrule\end{tabularx}\par\vspace{4mm}
\otIntertitre{14. Trigonométrie du triangle rectangle}\begin{tabularx}{\linewidth}{@{}Yl@{}}\toprule
Choisir cosinus, sinus ou tangente & \hyperref[exo-14-1]{14.1}, \hyperref[exo-14-2]{14.2}, \hyperref[exo-14-3]{14.3}, \hyperref[exo-14-4]{14.4}\\[2mm]
Calculer une longueur & \hyperref[exo-14-5]{14.5}, \hyperref[exo-14-6]{14.6}, \hyperref[exo-14-7]{14.7}, \hyperref[exo-14-13]{14.13}, \hyperref[exo-14-15]{14.15}\\[2mm]
Calculer un angle aigu & \hyperref[exo-14-8]{14.8}, \hyperref[exo-14-9]{14.9}, \hyperref[exo-14-10]{14.10}, \hyperref[exo-14-12]{14.12}\\[2mm]
Mesures indirectes et contrôle du modèle & \hyperref[exo-14-11]{14.11}, \hyperref[exo-14-12]{14.12}, \hyperref[exo-14-13]{14.13}, \hyperref[exo-14-14]{14.14}, \hyperref[exo-14-15]{14.15}\\[2mm]
\bottomrule\end{tabularx}\par\vspace{4mm}
\otIntertitre{15. Transformations et homothéties}\begin{tabularx}{\linewidth}{@{}Yl@{}}\toprule
Symétries, translations et rotations & \hyperref[exo-15-1]{15.1}, \hyperref[exo-15-2]{15.2}, \hyperref[exo-15-3]{15.3}, \hyperref[exo-15-4]{15.4}, \hyperref[exo-15-5]{15.5}\\[2mm]
Homothéties de rapports positifs et négatifs & \hyperref[exo-15-6]{15.6}, \hyperref[exo-15-7]{15.7}, \hyperref[exo-15-8]{15.8}, \hyperref[exo-15-13]{15.13}\\[2mm]
Facteurs de longueurs, aires et volumes & \hyperref[exo-15-9]{15.9}, \hyperref[exo-15-10]{15.10}, \hyperref[exo-15-14]{15.14}, \hyperref[exo-15-15]{15.15}\\[2mm]
Frises, pavages et rosaces & \hyperref[exo-15-11]{15.11}, \hyperref[exo-15-12]{15.12}\\[2mm]
\bottomrule\end{tabularx}\par\vspace{4mm}
\clearpage
\otPageTitre{Repères et exercices}{Chapitres 16 à 18}\fontsize{10}{13}\selectfont
\otIntertitre{16. Espace, solides et repérage sur la sphère}\begin{tabularx}{\linewidth}{@{}Yl@{}}\toprule
Perspectives, patrons, vues et sections & \hyperref[exo-16-6]{16.6}, \hyperref[exo-16-7]{16.7}, \hyperref[exo-16-8]{16.8}, \hyperref[exo-16-13]{16.13}\\[2mm]
Volumes usuels, boule et assemblages & \hyperref[exo-16-1]{16.1}, \hyperref[exo-16-2]{16.2}, \hyperref[exo-16-3]{16.3}, \hyperref[exo-16-4]{16.4}, \hyperref[exo-16-5]{16.5}, \hyperref[exo-16-11]{16.11}, \hyperref[exo-16-12]{16.12}, \hyperref[exo-16-14]{16.14}, \hyperref[exo-16-15]{16.15}\\[2mm]
Repérage dans un pavé & \hyperref[exo-16-9]{16.9}\\[2mm]
Latitude et longitude sur une sphère & \hyperref[exo-16-10]{16.10}\\[2mm]
\bottomrule\end{tabularx}\par\vspace{4mm}
\otIntertitre{17. Algorithmique, blocs et tableur}\begin{tabularx}{\linewidth}{@{}Yl@{}}\toprule
Variables, tests et exécution à la main & \hyperref[exo-17-1]{17.1}, \hyperref[exo-17-2]{17.2}, \hyperref[exo-17-4]{17.4}, \hyperref[exo-17-5]{17.5}, \hyperref[exo-17-6]{17.6}, \hyperref[exo-17-9]{17.9}, \hyperref[exo-17-10]{17.10}\\[2mm]
Boucles bornées et conditionnelles & \hyperref[exo-17-3]{17.3}, \hyperref[exo-17-7]{17.7}, \hyperref[exo-17-11]{17.11}\\[2mm]
Blocs, simulation et scripts parallèles & \hyperref[exo-17-8]{17.8}, \hyperref[exo-17-12]{17.12}, \hyperref[exo-17-13]{17.13}, \hyperref[exo-17-14]{17.14}\\[2mm]
Formules de tableur et recopie & \hyperref[exo-17-15]{17.15}\\[2mm]
\bottomrule\end{tabularx}\par\vspace{4mm}
\otIntertitre{18. Raisonner et préparer le brevet}\begin{tabularx}{\linewidth}{@{}Yl@{}}\toprule
Conjecture, preuve et contre-exemple & \hyperref[exo-18-1]{18.1}, \hyperref[exo-18-2]{18.2}, \hyperref[exo-18-6]{18.6}, \hyperref[exo-18-8]{18.8}, \hyperref[exo-18-10]{18.10}\\[2mm]
Hypothèses et contrôle des résultats & \hyperref[exo-18-3]{18.3}, \hyperref[exo-18-4]{18.4}, \hyperref[exo-18-5]{18.5}, \hyperref[exo-18-7]{18.7}, \hyperref[exo-18-9]{18.9}\\[2mm]
Problèmes de synthèse & \hyperref[exo-18-11]{18.11}, \hyperref[exo-18-12]{18.12}, \hyperref[exo-18-13]{18.13}, \hyperref[exo-18-14]{18.14}\\[2mm]
Méthode et format du brevet 2027 & \hyperref[exo-18-15]{18.15}\\[2mm]
\bottomrule\end{tabularx}\par\vspace{4mm}
\clearpage
